% ============================================
% DOC-26 (EN)
% English translation — SPEC-935-R355
% ============================================
\section*{\textbf{DOC-26} --- Oliveira's Notes --- The Mechanism}
\markboth{DOC-26 --- Oliveira's Notes --- The Mechanism}{}

\noindent \textit{Document extracted from the forensic file.}

\vspace{0.5em}\noindent\rule{\textwidth}{0.5pt}\vspace{0.5em}

\noindent {\footnotesize\texttt{PUBLIC ARCHIVE $\cdot$ DOCUMENTARY EVIDENCE $\cdot$ CASE 1,260}}\\[-0.2em]

\noindent\textbf{Document:} loose sheets found in the inner pocket of Dr. Heitor Oliveira's personal notebook, tied with twine. \textbf{Estimated date:} June 1979 to March 2026. \textbf{State:} the most recent sheets are written on different paper, in trembling, aged handwriting --- as if they had been written decades later, but tied to the same bundle.

\subsubsection*{Sheet 1 --- June 19, 1979}

\noindent Today I discovered what the diary does with time.

\noindent It is not just writing. It is \textit{addressing\textit{. When I write in Joaquim's diary, I can choose }when} that word will arrive. It is not conscious control --- it is like striking the right note and watching the glass vibrate. The page vibrates. And the word disappears.

\noindent I ran a test. I wrote: \textit{\textquotedblleft{}Maria, if you are reading this in 1980, light a candle in the window on August 15.\textquotedblright{}} I closed the diary. I put it under the bed. On August 15, 1979, I called the neighbor. She said: \textquotedblleft{}Your wife lit a candle in the window yesterday.\textquotedblright{}

\noindent It works.

\noindent The diary is not a record. It is a \textbf{rift}. And every rift has a signature.

\subsubsection*{\textit{Sheet 3 --- June 22, 1979}}

\noindent The pattern appeared today. I spent the night sending words to future dates and measuring the interval between writing and appearance.

\noindent The diary does not accept just any interval. It rejects it. The page stays blank, the ink runs, the pen scratches a dry stroke on its own. It only works in multiples of \textbf{62}.

\noindent 62 days. 124 days. 186 days. 248 days. 310 days. 372 days --- and so on.

\noindent Sixty-two. The number of years Joaquim spent at the Colony. The entity engraved the number in the diary as her signature. Every word that crosses time travels in leaps of 62. It is her breathing. The Cycle.

\noindent A letter destined for 2026 must travel in leaps of 62 days over 47 years. That is 276 leaps. Or 276 nights of writing. The page had already calculated before I finished the thought --- the number appeared written on its own in the right margin: \textbf{276}. She wants me to write. She wants the letter to arrive.

\subsubsection*{Sheet 7 --- illegible date (probably 1985)}

\noindent 2,190 days have passed. I write this note and hand it to the diary. Six years. Six complete cycles.

\noindent I look in the mirror and see age arriving. But the yellow eye does not age. The color grows stronger. Denser. More \textit{in charge}.

\noindent The diary showed me something new today: when I write a letter destined for a specific date, page 62 --- always page 62 --- lights up. That is where the rift opens. That is where the word leaps.

\noindent Page 62 of Joaquim's diary.

\noindent He wrote that page in 1919. Two years after entering the Colony. He had already discovered the number. He already knew that hunger counted in 62s.

\noindent He wrote: \textit{\textquotedblleft{}Sixty-two is her age. She told me today. Born in the drought of 1853, in the same trench, in the same smell. Sixty-two years later, I was born. Sixty-two years later, she found me.\textquotedblright{}}

\noindent The entity was born in 1853. Born of the drought. Of the hunger. Of the stacked bodies. She waits 62 years. Finds a host. Waits 62 more. Another host.

\noindent The Cycle is not human. It is geological. It is the time the earth takes to forget.

\subsubsection*{\textit{Sheet 12 --- March 2, 2003}}

\noindent Twenty-four years of cycles. 151 leaps of 62 days. The letter is traveling.

\noindent Sometimes I imagine the hand that will receive it. Someone in 2026. Someone who will open the archive and find an envelope with my name. They will think it is a prank. Then they will read the first line.

\noindent And then it will be too late.

\noindent The arriving letter is not a warning. It is a delivery. The rift opens on both sides. Whoever reads the letter in the future is also being read by the past.

\noindent I am sorry, Doctor. If you exist at all.

\subsubsection*{Sheet 15 --- last sheet, January 12, 2026}

\noindent The letter arrived.

\noindent I felt it. Today, January 12, 2026, I felt the rift close. The letter was received. Someone read it. The Cycle found the next one.

\noindent I look in the mirror and see an old man who is not me. Hanging skin. The right eye completely golden. The pupil --- I no longer know if it is mine. It elongates when the room darkens. It becomes vertical. A goat's slit. Of a creature that sees in the dark.

\noindent I delivered the diary today. I put it in a black cloth. I left it at the front desk. The woman at the desk asked my name. I said: \textit{\textquotedblleft{}Tell them it was patient 1,261.\textquotedblright{}}

\noindent She did not understand.

\noindent But the right person will understand.

\subsubsection*{\textit{Sheet 16 --- January 12, 2026 (night)}}

\noindent I came home after the delivery. I sat on the bed. I looked at my hands. 47 years. 47 years alive against my will. The letter I wrote to Maria in 1979 said I was going to protect others. That I was not going to let the diary be passed on. I believed that for decades. I believed I was in control.

\noindent I was not. I never was.

\noindent The hunger kept me alive. Not out of kindness. Not by chance. She needed a courier. Someone who could walk among the living, who had documents, identity, the appearance of a person. Someone who could walk into a public archive reception in 2026 and not raise suspicion. An 81-year-old man, in a shabby suit, with a package under his arm. Invisible. Perfect.

\noindent She planned this since the day I read the diary, in 1979. Since the day Joaquim wrote page 62, in 1919. Since the day the trench began to pulse, in 1853.

\noindent The Cycle did not ask me. The Cycle only used me.

\noindent I thought I was protecting Maria. I thought I was breaking the chain. But the chain never depended on me. I was only a link. And a link does not choose. A link holds. A link passes. Until it is another link's turn.

\noindent Maria died in 2004. I did not go to the funeral. Not because I did not want to. Because I \textit{could not}. On the day of the funeral, I woke up on a park bench in another city, with my clothes changed and 62 days of memory lost. She had taken me far away. She was not going to let anyone interrupt the delivery.

\noindent I am sorry, Maria. I am sorry, Doctor who will read this. I am sorry, next in line.

\noindent I had no choice.

\noindent And neither will you.

\noindent \textit{Archivist's note: this bundle of notes was found in February 2026, in a sealed brown envelope, inside the same box as patient 1,260's diary. The handwriting of the first sheets is firm, medical. The handwriting of the last ones is trembling, irregular, with strokes that resemble --- but are not identical to --- the calligraphy of Joaquim Antônio Correia. The date of the last note (January 12, 2026) coincides with the date the diary was returned to the Minas Gerais Public Archive for the first time.}
\noindent Page 62 of the original diary, where Joaquim wrote about the entity's origin in 1853, presents a moisture stain partially covering the text. The date \textquotedblleft{}1853\textquotedblright{} is written three times, in different handwritings --- as if it had been confirmed by different hands over time. The last confirmation is in a doctor's hand.

\vspace{0.5em}
\noindent\textit{\textbf{↪ Routes:} \rota{DOC-01} \rotasep \rota{LUC-02} \rotasep \rota{DOC-27}}
